\section{模型检验、误差分析与评价}\label{sec:evaluation}

\subsection{配对差异与条件不确定性}
对同一评价行，设候选$c$和参考$r_0$的平方误差差为
\begin{equation}
d_i=(y_i-\hat y_i^{(c)})^2-(y_i-\hat y_i^{(r_0)})^2,
\qquad\Delta_E=\frac1{n_E}\sum_{i\in E}d_i.
\label{eq:paired}
\end{equation}
$\Delta_E<0$表示候选MSE较小。既有区间通过有放回重采样相同行索引，取重采样均值的2.5\%与97.5\%分位。原六模型和顺序实验使用3000次，预处理对照使用2000次。预处理CV的两套差异先按同一行平均，再做行重采样，未将两套折当作独立的$2n_T$行。

表\ref{tab:intervals}给出核心比较。最终四模型相对原XGBoost的开发差异区间位于零以下，B区间则跨零；预处理改变的融合在两个集合区间均跨零。
\begin{table}[htbp]
\small
\setlength{\tabcolsep}{5pt}
\centering
\caption{固定预测下MSE差与95\%行配对区间}
\label{tab:intervals}
\begin{tabular}{L{0.08\linewidth}L{0.28\linewidth}L{0.17\linewidth}L{0.17\linewidth}L{0.17\linewidth}}
\toprule
集合 & 比较 & MSE差 & 下界 & 上界 \\
\midrule
开发 & 四模型−原XGB & -0.051565 & -0.101899 & -0.003706 \\
B & 四模型−原XGB & -0.018448 & -0.053797 & 0.016102 \\
开发 & 推荐处理−原处理 & -0.002787 & -0.011883 & 0.006530 \\
B & 推荐处理−原处理 & 0.000527 & -0.008715 & 0.009747 \\
\bottomrule
\end{tabular}
\end{table}


这些是给定预测和当前评价行下的条件区间，不包含重新训练、参数搜索、候选适应性、日期聚类与多重比较。开发集历史参与研究，故其区间不跨零也不能证明完整建模流程的无偏泛化收益。相对原大网络，原XGBoost的B误差差异区间也跨零，六模型MSE点估计排序不代表未来排序保证。

\subsection{目标分箱与尾部误差}
按真实血糖值划分$[0,6)$、$[6,10)$、$[10,+\infty)$，仅用于事后诊断，不用于预测分流或临床诊断。令$E_g$为分箱$g$内行集合，其平方误差贡献为
\begin{equation}
Q_g=\frac{\sum_{i\in E_g}(y_i-\hat y_i)^2}{\sum_{i\in E}(y_i-\hat y_i)^2}.
\label{eq:sse-share}
\end{equation}
平均预测减真实值用于描述偏差方向，负值表示该组平均低估。表\ref{tab:bands}给出最终融合的分箱结果。
\begin{table}[htbp]
\small
\setlength{\tabcolsep}{5pt}
\centering
\caption{最终四模型按真实目标分箱的误差诊断}
\label{tab:bands}
\begin{tabular}{lrrrrr}
\toprule
集合 & 目标区间 & 行数 & MSE & SSE占比 & 平均偏差 \\
\midrule
开发 & $[0,6)$ & 1598 & 0.449848 & 16.57\% & 0.341951 \\
开发 & $[6,10)$ & 336 & 1.923298 & 14.89\% & -0.931246 \\
开发 & $[10,+\infty)$ & 59 & 50.409513 & 68.54\% & -6.358480 \\
B & $[0,6)$ & 868 & 0.437047 & 30.73\% & 0.314046 \\
B & $[6,10)$ & 118 & 1.768811 & 16.91\% & -0.887488 \\
B & $[10,+\infty)$ & 14 & 46.158499 & 52.36\% & -6.396633 \\
\bottomrule
\end{tabular}
\end{table}


开发集中59条真实值至少10的记录贡献约68.54\%平方误差，B集中14条贡献约52.36\%，平均偏差分别约$-6.358480$和$-6.396633$。较少高值记录对总体误差有较大影响，最终融合未消除尾部低估。

总体改善也不表示每一分箱都改善。B中$[6,10)$组，原XGBoost的MSE为1.748310，最终融合为1.768811，略有上升。上述分箱不足以识别尾部误差的病因或测量机制，不能通过删除这些高损失记录制造更低MSE。

\subsection{竞赛目标的证据覆盖}
在精度方面，研究保留相同评价行，通过基准、阶段和固定预处理对照提供可比结果；四模型在当前开发/B的MSE点估计均低于原XGBoost。但B区间与历史使用限制了稳定优势的解释。

在计算方面，研究给定决策树24、AdaBoost18、原XGBoost24及随机森林12个参数候选；后续联合搜索为100/12/8，融合为有限权重网格。预算使计算过程可界定，但不等于已测得四模型训练或推理速度优于其他模型。不同实验的存档耗时也未形成统一硬件、相同任务的横向基准。

在解释性方面，OLS/Lasso用于理解给定线性表示，特征组消融用于观察受控候选变化，分箱用于定位误差贡献。最终融合构成和权重明确，但尚未实施个体解释、因果分析或临床验证。因此，本文不将探索系数解释为最终融合的个体医学贡献。

\subsection{局限与后续研究方向}
两套训练CV均用于选择，候选已在第一套折筛选，网络还包含训练诊断驱动的补充，故不是完整嵌套验证。开发与B的历史使用也不能通过冻结本次配置逆转。进一步评估需要新的独立样本，或围绕完整选择过程重新设计验证，不能继续依已知B最低分选择模型并称其为最终盲测。

现有ID无法排除同人多次体检，日期范围重叠，患者级与日期级依赖未在行区间中充分体现。单位、采样条件和原始标签排序来源缺项，也限制医学解释。网络未开展额外初始化种子重复，两个CV种子主要描述划分变化，不覆盖全部训练随机性。

预处理对照固定已有参数，结果不代表每种变换重新调参后的最优比较。本文未实施系统超参数扰动、当前四模型严格时间外推或最终融合个体解释实验，故不将其写成已完成的敏感性和临床有效性检验。未来可优先核验尾部记录的输入与测量背景，并在新数据上检验收益，尚不能从现有结果确定增加何种输入即可消除误差。
